\documentclass[sigconf,nonacm]{acmart}
\usepackage{booktabs}
\usepackage{amsmath}
\usepackage{graphicx}
\usepackage{xcolor}
\usepackage{array}
\setlength{\emergencystretch}{2em}

% ======================================================================
% WORKING FINAL DRAFT / 24 NOVEMBER 2026
% Figures and table cells marked PENDING are NOT observed results.
% Keep the completed Run A measurements separate from later experiments.
% Search for "FINAL UPDATE" to locate every required replacement.
% Final submission: check 8-page ACM limit, authors, references, and figures.
% ======================================================================
\definecolor{pendinggray}{gray}{0.42}
\newcommand{\pending}[1]{\textcolor{pendinggray}{\emph{[PENDING: #1]}}}
\newcommand{\tbd}{\textcolor{pendinggray}{\textit{TBD}}}
\newcommand{\figslot}[2]{%
  \fbox{\begin{minipage}[c][#1][c]{0.92\columnwidth}
  \centering\small\textcolor{pendinggray}{\textbf{[FIGURE PENDING]}}\\[5pt]
  \textcolor{pendinggray}{#2}
  \end{minipage}}}
\newcommand{\widefigslot}[2]{%
  \fbox{\begin{minipage}[c][#1][c]{0.92\textwidth}
  \centering\small\textcolor{pendinggray}{\textbf{[FIGURE PENDING]}}\\[5pt]
  \textcolor{pendinggray}{#2}
  \end{minipage}}}

\title[Counterfactual Traffic Accident Spillovers]{Causal Graph Learning for Counterfactual Estimation of Traffic Accident Spillovers}
% Author block preserved from the original proposal. Confirm authorship
% and contribution language before final submission.
\author{Kyle Luong}
\affiliation{\institution{University of Virginia}\city{Charlottesville}\state{Virginia}\country{USA}}
\author{Nate Mackintosh}
\affiliation{\institution{University of Virginia}\city{Charlottesville}\state{Virginia}\country{USA}}
\author{Maddox Shrum}
\affiliation{\institution{University of Virginia}\city{Charlottesville}\state{Virginia}\country{USA}}
\renewcommand{\shortauthors}{Luong et al.}

\begin{document}

\begin{abstract}
Traffic crashes may change speeds on the affected road and on roads
connected to it. Distinguishing a crash's contribution from ordinary
congestion is difficult because crashes are not randomly located in
space or time and neighboring roads share other sources of variation.
We study this problem using Chicago crash records, ten-minute arterial
speed measurements, and a road-segment graph. Our evaluation compares
matched event studies, conventional and graph-based forecasts of
untreated speed, and treatment-aware models that predict outcomes under
specified crash exposures. The design tests three distinct questions:
how estimated effects vary across network distance and time; whether
real road connectivity improves prediction and recovery of known
effects; and whether observational neighboring-road contrasts withstand
placebo and sensitivity checks. A completed 2024--2026 run establishes
the data pipeline and initial baselines. Controlled graph comparisons,
expanded semi-synthetic interventions, and improved identification
diagnostics supply the final tests. We report forecasting, known-effect
recovery, and real-event estimates separately so that conclusions need
not favor any particular graph model or assume a detectable spillover.
\end{abstract}
\keywords{graph neural networks, counterfactual traffic prediction,
network interference, traffic crashes, causal inference, spatiotemporal graphs}
\maketitle

\section{Introduction}
A traffic crash can affect the road where it is reported and may alter
conditions elsewhere through queues, rerouting, or shared congestion.
Traffic speed after a crash, however, is not the same as the speed
\emph{caused} by the crash. Congestion can precede an incident, demand
shocks can affect many roads simultaneously, and available speed
measurements may change across road segments and times. The causal
question compares the speed observed after a reported crash with the
speed that would have occurred on that road at that time without it.

Urban roads form a network, making graph machine learning (GML) a
natural modeling tool. Yet the value of the graph must be tested, not
assumed. A graph neural network (GNN) may improve ordinary speed
forecasts without improving estimates of crash effects. Moreover,
observed network adjacency need not capture the actual path of traffic
propagation. We therefore compare graph and non-graph models while
separately evaluating the credibility of exposure and control choices.

Our study preserves the original proposal's focus on \emph{effect
estimation rather than prediction alone}. We organize the experiments
around three research questions:
\begin{itemize}
  \item \textbf{RQ1: Effect patterns.} What speed changes are associated
  with crashes on the incident road and at one and two graph hops over
  the next hour? Do these contrasts survive pre-event, placebo,
  overlapping-crash, and observation-coverage diagnostics, and how
  long does a detectable contrast last?
  \item \textbf{RQ2: Graph contribution.} Does the real road graph
  improve held-out forecasting and \emph{known} effect recovery beyond
  equally informed non-graph and false-topology comparisons?
  \item \textbf{RQ3: Explicit treatment modeling.} Does a
  treatment-conditioned GNN estimate intervention responses more
  accurately than a treatment-conditioned non-graph model or the
  existing matched-residual estimator?
\end{itemize}

The contributions are an empirical framework combining observational
and known-effect evaluation on an urban road graph; a controlled test
of graph information versus model capacity and input access; and a
transparent account of what pre-event diagnostics permit us to infer
about reported crashes. We do not presume which model or exposure
pattern will perform best.

\section{Related Work}
\textbf{Crash impact estimation.} Cao et al.\ quantify incident
impacts using spatiotemporal regions and comparison traffic
\cite{cao2021traffic}. Li et al.\ estimate heterogeneous highway
crash effects using doubly robust causal learning
\cite{li2024crash}. Both motivate comparisons against plausible
untreated conditions rather than raw post-crash speeds.

\textbf{Graph traffic forecasting.} GCN and GraphSAGE provide general
message-passing encoders \cite{kipf2017gcn,hamilton2017graphsage},
while diffusion convolutional recurrent networks demonstrate the
value of road graphs for time-series prediction
\cite{li2018dcrnn}. Incident-aware traffic models such as CASTANET
also study disruption-sensitive forecasting
\cite{kitahara2026castanet}. These are relevant predictive baselines,
but their speed error alone does not evaluate causal-effect accuracy.

\textbf{Causal learning under interference.} Network Deconfounder
learns graph representations for individual causal inference
\cite{guo2020network}; Ma and Tresp study interventions whose effects
can propagate over networks \cite{ma2021interference}. Observational
network adjustment requires explicit assumptions on exposure and
confounding \cite{forastiere2022interference}. MSCT addresses
treatment-aware post-crash counterfactual prediction with
time-varying confounding \cite{li2024msct}. The act of adding treatment
heads is therefore not itself our novelty claim. Recent research
also questions whether fixed hop-based exposure mappings are adequate
\cite{huber2026exposure}. Our emphasis is a shared test of topology,
treatment conditioning, and falsification in the same urban study.

\section{Data and Causal Targets}
\subsection{Chicago Road Observations}
Chicago provides geocoded crash reports and historical arterial speed
estimates \cite{chicagoCrashes,chicagoSpeeds}. We use the half-open Run
A window June 11, 2024--May 1, 2026. Its converted inputs include
101,138,608 speed records and 207,612 reported crashes. The matched
monitored network has 881 segments and 4,744 directed arcs for graph
modeling. Only 57.1\% of the road-by-ten-minute grid has observed speed
values in the primary joined data. This is a study of roads covered by
the source, not every Chicago crash or street.

The raw speed feed also contains explicit unavailable-speed readings;
these are distinct from absent road/time cells. We retain observation
masks and avoid substituting missing speed with a measured zero.
Because the semantics of recorded zero-speed values are uncertain, a
separate sensitivity run excludes them and reports the changed sample
rather than treating results as directly paired.

\begin{table}[t]
\centering
\caption{Completed Run A data selection. Later event-study tables
use their explicitly stated event and outcome eligibility rules.}
\label{tab:funnel}
\footnotesize
\begin{tabular}{lr}
\toprule
Selection stage & Number \\
\midrule
Crashes reported in study period & 207,612 \\
With usable coordinates & 205,387 \\
Matched to monitored segment & 44,108 \\
On a segment with speed history & 44,019 \\
Primary matched-analysis events & 32,479 \\
Common 0/1/2-hop event-study cohort & 10,162 \\
\bottomrule
\end{tabular}
\end{table}

Crash-to-road matching and event eligibility reduce the represented
population substantially (Table~\ref{tab:funnel}). The primary
matched-residual set and the broader event-study processing set use
slightly different rules; horizon-specific sample sizes depend on
observed outcomes and eligible controls. These differences are
reported, not silently pooled. Frozen road-node identifiers are
retained across all graph ablations and treatment definitions.

\subsection{Exposure, Potential Outcomes, and Assumptions}
Let $G=(V,E)$ be the observed road-segment graph. Crash event $e$ occurs
at segment $s_e$ and recorded time $t_e$. For evaluated road $i$,
let $A_{ei}\in\{\varnothing,0,1,2\}$ mean no qualifying exposure, a
crash on the road itself, exposure one hop away, or exposure two hops
away. An overlap rule excludes or resolves multiple nearby incidents.
Distance is computed on the \emph{real} road graph and is kept fixed
when testing alternative model graphs.

For horizon $h\in\{10,20,30,40,50,60\}$ minutes, define
$Y_{eih}(a)$ as speed under exposure $a$ and
$Y_{eih}(\varnothing)$ as speed without that specified exposure.
The target contrast is
\begin{equation}
 \tau_{eih}(a)=Y_{eih}(a)-Y_{eih}(\varnothing).
 \label{eq:effect}
\end{equation}
Negative contrasts indicate lower speeds under exposure. Only the
factual outcome is observed for a real event. The observed graph and
exposure rule are approximations: vehicle routing depends on direction,
turns, road capacity, and unmonitored links. Identification also
requires sufficiently comparable untreated conditions and appropriate
handling of concurrent crashes, hidden demand changes, and reporting
times. These conditions are examined, not assumed proven.

\begin{figure}[t]
\centering
\figslot{0.84in}{Automated network schematic: crash road (0-hop),\newline first and second neighboring rings; distinguish\newline graph exposure from physical propagation.}
\Description{Placeholder for a schematic displaying a crash road and first and second neighboring exposure rings.}
\caption{\textbf{Planned study diagram.} Illustrate the exposure
categories and factual versus no-crash outcome comparison. This is a
schematic, not evidence that traffic physically follows each edge.}
\label{fig:design}
\end{figure}

\section{Methods}
\subsection{Matched Event Study and Residual Estimates}
We construct crash and control windows using the same monitored roads,
clock/calendar context, and traffic recorded \emph{before} the event.
For crash event $e$, target road $i$, pre-event reference time $b$,
and matched untreated control window $c$, an event-study contrast is
\begin{equation}
 \widehat\tau^{ES}_{eic,h}
 =(Y_{eih}-Y_{eib})-(Y_{cih}-Y_{cib}).
 \label{eq:did}
\end{equation}
The final ring-by-horizon estimate averages matched controls within an
event and observed target roads within a ring before aggregating
events. This prevents events with many monitored neighbors from
receiving automatic extra weight. We also report estimates based on
all available observations and on a fixed event cohort present at
all hops and horizons.

Separately, a model $f$ predicts the expected speed in an untreated
window from pre-event histories. Let
$r_{eih}=Y_{eih}-\widehat Y^{(0)}_{eih}$ be the observed minus
no-crash-predicted speed. The residual estimator compares the average
crash residual to that of matched unexposed periods:
\begin{equation}
 \widehat\tau^{res}_{k,h}
 =\overline r_{\mathrm{crash},k,h}
  -\overline r_{\mathrm{control},k,h}.
 \label{eq:residual}
\end{equation}
This can reduce predictable forecast bias, but matching and residual
adjustment alone cannot remove hidden confounding.

\subsection{Untreated Speed Forecasters}
The current pipeline uses twelve ten-minute pre-event bins to predict
six ten-minute horizons. It evaluates persistence, historical mean,
MLP, temporal GRU, GCN, GraphSAGE, and STGCN on masked observed
targets. Model inputs are histories and road information available
before the evaluated event; post-event speed is never used to
construct a hypothetical no-crash input. Chronological splits,
window-overlap exclusion, and untreated training targets are checked
against the frozen evaluation specification.

\subsection{Treatment-Aware Graph and Non-Graph Models}
A separate model explicitly conditions on crash exposure. It is
implemented and queued; no result from it is reported here.
An encoder $z_{ei}=f_\theta(X_{e,\le t_e},G)$ uses only pre-event
information, followed by an exposure-conditional outcome function:
\begin{equation}
 \widehat Y_{eih}(a)=g_\phi(z_{ei},a,h),\qquad
 \widehat\tau_{eih}(a)=\widehat Y_{eih}(a)
                      -\widehat Y_{eih}(\varnothing).
 \label{eq:heads}
\end{equation}
The training loss is applied \emph{only} to the outcome corresponding
to each observed factual exposure; no invented counterfactual targets
are used for real events. We compare a non-graph MLP with a GCN- or
STGCN-based encoder. Importantly, the MLP receives the same
pre-event road signals and explicit summaries of neighboring speeds,
so a graph improvement cannot be attributed only to access to neighbor
information. Exposure labels, outcomes, splits, masking, and training
budget are held fixed. Potential-outcome prediction is a modeling
approach, not proof of unconfounded treatment assignment.

\section{Experimental Evaluation}
\subsection{Shared Splits and Fair Graph Comparisons}
Our controlled experiment changes only model message-passing topology:
real road edges, self-only/no neighboring messages, and one documented
rewired graph. Treatment labels, target roads, raw input observations,
chronological split, masks, and the underlying physical graph remain
unchanged. The graph encoder and comparably informed non-graph model
are trained and scored on the same eligible instances. We report
forecast MAE/RMSE separately from known-effect MAE/bias, with at least
three model seeds when resource constraints permit. Event bootstrap
uncertainty and across-seed variation answer different questions and
are not conflated.

\subsection{Known-Effect Semi-Synthetic Evaluation}
We programmatically inject interventions into held-out crash-free
traffic histories and preserve both untreated and treated potential
outcomes, exposure assignments, ground-truth contrasts, event IDs,
and missingness masks. The initial 200-event additive benchmark obeys
$Y^{sim}=Y^{(0)}+\delta$. For a residual-style estimator,
\begin{equation}
 (Y^{sim}-\widehat Y^{(0)})-\delta
  =Y^{(0)}-\widehat Y^{(0)}.
 \label{eq:identity}
\end{equation}
Thus individual injected-effect error can be exactly ordinary
untreated forecasting error. We retain this test for continuity, but
it is not standalone evidence of causal identification.

The extended benchmark separately tests (i) a direct effect with
\emph{zero} true 1/2-hop spillover, (ii) known time-varying spatial
decay with multiple magnitudes, and (iii), if validated, treatment
assignment dependent on pre-event congestion rather than random
placement. Models are evaluated on the same held-out interventions
using hop-/horizon-specific effect MAE/RMSE, average-effect bias,
false-positive neighboring effects, and paired event-level intervals.
A negative or null result is informative if supported by these tests.

\subsection{Observational Diagnostics and Conditional Extensions}
Our strongest observational checks are target-road pre-crash balance,
pre-event pseudo-crash times, concurrent-crash flagging on the target road
and its neighbors, and common-event trajectories. We also
compare observation availability by exposure and horizon, because
complete-case estimates could change when sensors or source coverage
are correlated with incidents. Recorded zero speeds receive a
separate sensitivity analysis.

We also summarize how long a detectable contrast lasts in each ring,
counting horizon steps whose event-clustered interval lies wholly below
zero. It reuses the event-study estimates rather than testing anything
new, and is read only after the placebo, matching and missingness
diagnostics, because a duration taken from intervals the placebo cannot
clear describes the arms rather than the crash. A ring still negative at
the final step is reported as censored: the observation window ends
before the effect does.

\section{Results: Completed Baseline and Final Tests}
\subsection{Completed Run A: Model and Data Checks}
Table~\ref{tab:forecast} reports the first completed Run A held-out
forecast comparison. Its shared test split contains 14,754 forecast
windows. These observations are a reproducible starting point for the
controlled graph and treatment-aware experiments, \emph{not} a final
ranking of counterfactual estimators.

\begin{table}[t]
\centering
\caption{Initial Run A held-out speed prediction, in mph. These
numbers are complete for the original seven-model run; subsequent
topology-controlled comparisons are reported separately.}
\label{tab:forecast}
\footnotesize
\begin{tabular}{lrr}
\toprule
Method & Test MAE & Test RMSE \\
\midrule
Persistence & 4.141 & 5.623 \\
Historical mean & 3.151 & 4.266 \\
MLP & 3.102 & 4.231 \\
Temporal GRU & 3.067 & 4.206 \\
GCN & 3.072 & 4.197 \\
GraphSAGE & 3.036 & 4.165 \\
STGCN & 3.029 & 4.165 \\
\bottomrule
\end{tabular}
\end{table}

The original additive synthetic benchmark produced pooled effect MAEs
of 3.247 mph for historical mean and 3.288 mph for STGCN, with other
model scores recorded in the experiment logs. Those results are
reported descriptively, subject to Eq.~\ref{eq:identity}; the final
known-effect comparison will use the expanded interventions below.

\subsection{Crash-Road and Neighbor Contrasts}
The initial available-case event study at 30 minutes reports
matched contrasts of $-0.523$, $-0.106$, and $-0.056$ mph for the
incident road, 1-hop roads, and 2-hop roads. At 10 minutes, the 1-hop
contrast is $-0.098$ mph, while its corresponding placebo is
$-0.069$ mph. These are \emph{observational contrasts}, not verified
causal spillovers. The common cohort (10,162 events) and revised
neighbor-specific controls form part of the final sensitivity analysis.

\begin{figure*}[t]
\centering
\widefigslot{1.12in}{Panel A: 0/1/2-hop contrasts with event-level 95\% intervals\\[3pt]
Panel B: fixed-cohort estimates, pre-event placebo and neighbor-specific controls}
\Description{Placeholder for two panels of observed crash and neighboring-road contrasts across the first hour, including confidence intervals and placebo comparisons.}
\caption{\textbf{Final observational effect plot (pending).} Show 10--60-minute
trajectories, available and fixed cohorts, signed placebo comparisons,
matching version, and event counts. Source: validated
\texttt{outputs/identification-runA/}.}
\label{fig:eventcurves}
\end{figure*}

\begin{table}[t]
\centering
\caption{\textbf{Observational checks (pending).} Supply
pre-event imbalance, placebo magnitude, concurrent-crash sensitivity,
and outcome availability; do not interpret placebo share as a p-value.}
\label{tab:diagnostics}
\footnotesize
\begin{tabular}{p{0.43\columnwidth}p{0.19\columnwidth}p{0.19\columnwidth}}
\toprule
Diagnostic & Before & Revised \\
\midrule
Neighbor pretrend gap & \tbd & \tbd \\
1-hop pre-event placebo (mph) & \tbd & \tbd \\
Concurrent-crash-excluded contrast & \tbd & \tbd \\
Exposed/control speed availability & \tbd & \tbd \\
Fixed-cohort 1-hop effect & \tbd & \tbd \\
\bottomrule
\end{tabular}
\end{table}

% FINAL UPDATE: Replace this sentence with measured pretrend/placebo,
% missingness, concurrent-crash and common-event findings, not a favored claim.
\pending{Interpret Figure~\ref{fig:eventcurves} and Table~\ref{tab:diagnostics}
using the finalized exposure rule and matched event counts.}

\subsection{Graph Topology and Treatment-Aware Models}
The decisive model comparisons hold target data fixed while changing
topology or treatment awareness. Table~\ref{tab:causalmodels} and
Figure~\ref{fig:ablation} are reserved for the final verified
results. Only synthetic experiments with known ground truth permit
direct effect-error comparisons.

\begin{table}[t]
\centering
\caption{\textbf{Paired model comparisons (pending).}
Every row uses a declared common evaluation set and treatment labels
from the real graph. Lower error is better.}
\label{tab:causalmodels}
\footnotesize
\begin{tabular}{lccc}
\toprule
Model / graph & Speed MAE & Effect MAE & Bias \\
\midrule
Encoder / real edges & \tbd & \tbd & \tbd \\
Encoder / self-only & \tbd & \tbd & \tbd \\
Encoder / rewired & \tbd & \tbd & \tbd \\
Treatment MLP + neighbor inputs & \tbd & \tbd & \tbd \\
Treatment GNN / real edges & \tbd & \tbd & \tbd \\
\bottomrule
\end{tabular}
\end{table}

\begin{figure}[t]
\centering
\figslot{0.98in}{Two aligned panels: (a) held-out forecasting error;\\
(b) known-effect recovery error; real/self-only/rewired topology.}
\Description{Placeholder for paired panels comparing forecasting error and synthetic treatment-effect error across graph ablations.}
\caption{\textbf{Graph contribution (pending).} Plot paired
prediction and treatment-effect errors on consistent observations,
with training-seed spread where available.}
\label{fig:ablation}
\end{figure}

% FINAL UPDATE: Insert verified model comparison and explain, regardless
% of which model has the lowest effect MAE.
\pending{Report paired differences and uncertainty for graph and
non-graph counterfactual models, separately from speed forecasting.}

\subsection{Synthetic Null and Propagation Checks}
\begin{table}[t]
\centering
\caption{\textbf{Known-effect benchmarks (pending).} Compare estimated
versus true effects using the same held-out intervention IDs.}
\label{tab:synthetic}
\footnotesize
\begin{tabular}{p{0.43\columnwidth}cc}
\toprule
Intervention & Primary outcome & Result \\
\midrule
Direct only; true neighbor effect $=0$ & False 1/2-hop effect & \tbd \\
Known spatial/temporal decay & Hop-specific effect MAE & \tbd \\
Pre-congestion-dependent assignment & Effect bias & \tbd \\
\bottomrule
\end{tabular}
\end{table}

% FINAL UPDATE: Insert null false-positive rates, decay recovery, and
% treatment-assignment sensitivity, with known effect sizes and seeds.
\pending{Insert a final synthetic result figure if it offers evidence
not already conveyed by Table~\ref{tab:synthetic}.}

\section{Discussion and Limitations}
The main scientific interpretations depend on results that are still
being evaluated. A model with lower speed error may or may not recover
known crash effects more accurately. If real edges outperform
self-only and rewired graph variants on identical treatment-effect
targets, that supports a contribution from road connectivity under
the tested data-generating conditions. If the result is limited to
forecasting, the conclusion should be limited accordingly. If
explicit treatment conditioning changes effect accuracy, the
comparison with an equally supervised MLP is needed to determine
whether the gain is specifically graph-related.

Observational differences require more caution. A slowdown estimated
for a nearby road can reflect a crash, earlier congestion, another
incident, route diversion, an incorrect exposure map, or a change in
which speed readings are available. Statistical intervals account for
specified sampling uncertainty but do not remove bias from unobserved
confounding. In addition, the average change in mph may conceal
important variation by pre-crash traffic and incident characteristics.
Direction and heterogeneity are therefore interpreted only where the
recorded network and data coverage support them.

We also considered an unexposed-road negative control, comparing rings
beyond the two-hop exposure radius. We do not report one. Past that radius the
crash arm stops being exposure-masked, so the control arm's clean-post filter
begins applying to both arms alike, and the resulting step at that distance
would measure the filter rather than traffic. We state this as a limitation
instead of reporting a contrast we cannot interpret.

We do not separate upstream from downstream neighbors, and we report no
subgroup analysis. The segment graph links two segments whenever their
underlying road edges meet at an intersection, so it is undirected by
construction and carries no flow orientation. Segment direction labels do
exist and would orient roughly thirty percent of one-hop pairs, but
exploiting them would require rebuilding the exposure graph, and the
remaining pairs would still depend on turn restrictions the data does not
record. We leave directional and heterogeneous effects as future work
rather than report a partial version of either.

The study is limited to matched incidents on monitored Chicago arterial
segments, uses reported crash times and locations, and approximates
interference using a finite observed graph. The synthetic benchmarks
make treatment effects known \emph{by construction}, which is valuable
for model evaluation but does not prove identification for real
crashes. The final discussion will distinguish a detected association,
a model's recovery under known intervention, and a causal effect
supported by explicit assumptions.

\section{Conclusion and Reproducibility}
We provide an evaluation framework for traffic crash effects that
links a real urban road graph, matched observational comparisons,
non-graph and graph-based prediction, and interventions with known
ground truth. The completed Run A establishes feasibility at city
scale; the final experiments test the contribution of actual road
connections and explicit crash exposures while examining whether
neighbor-road comparisons survive falsification diagnostics.

% FINAL UPDATE: Replace only after all tables are frozen. Include one
% neutral sentence each for RQ1/RQ2/RQ3, supported by final evidence.
\pending{Insert concise final findings for RQ1--RQ3, including null,
mixed, and inconclusive results where supported.}

For reproducibility, the experiments are queued by a single command that
submits a dependency-ordered job graph; each stage records a completion marker,
so a resubmission resumes rather than repeats. Every job stores its commit,
working-tree diff, source archive, package lock, Slurm settings and seeds. A
generated result index names the file supplying each manuscript table and
figure, and marks anything unfinished as missing rather than filling it. Existing completed runs are preserved separately from each
follow-on experiment. Code and scripted analyses automatically
construct labels, event cohorts, figures, and result tables; the
counterfactual outcomes for real crashes remain unobserved.

\bibliographystyle{ACM-Reference-Format}
\bibliography{references}
\end{document}
